\documentclass[aspectratio=169,11pt]{beamer}

\usepackage[spanish,es-nodecimaldot]{babel}
\usepackage[T1]{fontenc}
\usepackage[utf8]{inputenc}
\usepackage{xcolor}
\usepackage{booktabs}
\usepackage{listings}
\usepackage{pgfplots}
\usepackage{tikz}
\usetikzlibrary{positioning,arrows.meta}
\pgfplotsset{compat=1.18}

\definecolor{azul}{HTML}{174A7E}
\definecolor{azulclaro}{HTML}{DCECF8}
\definecolor{verde}{HTML}{2E7D32}
\definecolor{gris}{HTML}{F3F5F7}
\definecolor{naranja}{HTML}{A64B00}
\definecolor{rojo}{HTML}{B71C1C}

\usetheme{Madrid}
\setbeamercolor{structure}{fg=azul}
\setbeamercolor{frametitle}{fg=azul}
\setbeamercolor{block title}{bg=azul,fg=white}
\setbeamercolor{block body}{bg=azulclaro,fg=black!85}
\setbeamercolor{alerted text}{fg=rojo}
\setbeamertemplate{navigation symbols}{}
\setbeamertemplate{footline}{\hfill\insertframenumber\hspace{0.4cm}\vspace{0.2cm}}

\lstdefinestyle{codigo}{
  language=Python,
  basicstyle=\ttfamily\scriptsize,
  backgroundcolor=\color{gris},
  frame=single,
  breaklines=true,
  keepspaces=true,
  showstringspaces=false,
  columns=fullflexible,
  keywordstyle=\color{azul}\bfseries,
  commentstyle=\color{verde},
  stringstyle=\color{naranja}
}

\title{El Pipeline Completo de IA}
\subtitle{Banderazo inicial del proyecto final (10 de noviembre -- 15 de enero)}
\author{Laboratorio de Inteligencia Artificial}
\date{}

\begin{document}

\begin{frame}
  \titlepage
\end{frame}

\begin{frame}{El objetivo de hoy: unir todas las piezas}
  \begin{center}
    \Large De ejercicios aislados a un \textbf{sistema completo y verificable}
  \end{center}
  \vfill
  \begin{itemize}
    \item Durante el semestre construimos ladrillos individuales: agentes, redes neuronales, espacios latentes, vectores semánticos y auditoría de tablas CSV.
    \item Un modelo aislado no es un proyecto: un proyecto es un \textbf{pipeline reproducible} desde el dato crudo hasta la autopsia de sus errores.
    \item Hoy damos el banderazo de salida: tenemos \textbf{dos meses} (cierre: \textbf{15 de enero}) para diseñar, construir y defender un pipeline completo.
  \end{itemize}
\end{frame}

\begin{frame}{Mapa de lo que ya sabemos hacer}
  \begin{columns}[T]
    \begin{column}{0.32\textwidth}
      \begin{block}{1. Datos y Supervisión}
        \small
        \begin{itemize}
          \item Auditoría de CSV y limpieza.
          \item Anomalías (Benford, IQR).
          \item Regresión y Redes Neuronales.
          \item \emph{Early stopping} e hiperparámetros.
        \end{itemize}
      \end{block}
    \end{column}
    \begin{column}{0.32\textwidth}
      \begin{block}{2. Espacio Latente}
        \small
        \begin{itemize}
          \item Emergencia y autoorganización.
          \item Mapas de Kohonen (SOM).
          \item Autoencoders clásicos.
          \item Modelos generativos (VAE).
        \end{itemize}
      \end{block}
    \end{column}
    \begin{column}{0.32\textwidth}
      \begin{block}{3. Texto y Contexto}
        \small
        \begin{itemize}
          \item Embeddings semánticos.
          \item Similitud coseno y $k$-NN.
          \item Fragmentación y RAG.
          \item Auto-atención en Transformers.
        \end{itemize}
      \end{block}
    \end{column}
  \end{columns}
  \vfill
  \centering
  \textbf{Regla de oro del proyecto:} no hay que meter todos los modelos a la fuerza; hay que elegir la ruta adecuada para una pregunta clara.
\end{frame}

\begin{frame}{Las 6 etapas de un Pipeline de IA}
  \centering
  \begin{tikzpicture}[
    node distance=0.35cm,
    caja/.style={draw=azul, fill=azulclaro, rounded corners=3pt, minimum width=2.2cm, minimum height=1.15cm, align=center, font=\scriptsize\bfseries}
  ]
    \node[caja] (e1) {1. Pregunta\\y Auditoría};
    \node[caja, right=of e1] (e2) {2. Limpieza y\\Transformación};
    \node[caja, right=of e2] (e3) {3. Revisión de\\Atípicos};
    \node[caja, right=of e3] (e4) {4. Partición\\sin Fuga};
    \node[caja, right=of e4] (e5) {5. Referencia\\vs. Modelo};
    \node[caja, right=of e5, fill=verde!20, draw=verde] (e6) {6. Autopsia\\y Límites};

    \draw[-{Stealth[length=2.5mm]}, thick, azul] (e1) -- (e2);
    \draw[-{Stealth[length=2.5mm]}, thick, azul] (e2) -- (e3);
    \draw[-{Stealth[length=2.5mm]}, thick, azul] (e3) -- (e4);
    \draw[-{Stealth[length=2.5mm]}, thick, azul] (e4) -- (e5);
    \draw[-{Stealth[length=2.5mm]}, thick, verde] (e5) -- (e6);
  \end{tikzpicture}
  \vfill
  \begin{itemize}
    \item Si falla la Etapa 1 o 2, el modelo aprenderá basura (\emph{garbage in, garbage out}).
    \item Si falla la Etapa 4, el resultado será un espejismo por fuga de información.
    \item Si falta la Etapa 5 o 6, no sabremos si el modelo aporta valor real o dónde se rompe.
  \end{itemize}
\end{frame}

\begin{frame}{Tres rutas posibles para su proyecto final}
  \begin{description}
    \item[Ruta A — Tabular y Predicción Supervisada:]
      CSV propio o de simulación (energía, demanda, sensores del robot). Limpieza documentada, detección de atípicos con IQR, regresión o red neuronal con \emph{early stopping}, superando a una referencia en datos de prueba.
    \vspace{0.2cm}
    \item[Ruta B — Representación y Espacio Latente:]
      SOM, Autoencoder o VAE sobre datos multivariados o registros de episodios para descubrir agrupamientos, comprimir dimensiones, interpolar estados o detectar anomalías de reconstrucción.
    \vspace{0.2cm}
    \item[Ruta C — Texto, Embeddings y Sistema RAG:]
      Corpus de documentos reales del equipo, comparación de estrategias de \emph{chunking}, base vectorial con similitud coseno y evaluación rigurosa de aciertos y fallos de recuperación.
  \end{description}
\end{frame}

\begin{frame}[fragile]{Etapas 1 a 3 en código: auditar antes de modelar}
  \begin{lstlisting}[style=codigo]
def limpiar_y_transformar(filas_crudas):
    limpios = []
    for fila in filas_crudas:
        try:
            equipos = float(fila["equipos_activos"])
            consumo = float(fila["consumo_kwh"])
        except ValueError:
            continue  # Celda vacia o texto invalido
        if equipos < 0 or consumo <= 0:
            continue  # Regla de dominio: fisicamente imposible
        limpios.append(fila)
    return limpios
  \end{lstlisting}
  \begin{block}{Distinción clave para su proyecto}
    Un valor \textbf{imposible} (ej. $-5$ equipos activos) se corrige o descarta por regla de dominio. Un valor \textbf{extremo pero posible} (marcado por IQR) se investiga antes de decidir si se excluye.
  \end{block}
\end{frame}

\begin{frame}{Etapa 4: la trampa mortal de la fuga de datos}
  \begin{columns}[T]
    \begin{column}{0.48\textwidth}
      \begin{block}{Fuga temporal (incorrecto)}
        \small
        Mezclar al azar días de diciembre en el entrenamiento y evaluar en días de octubre.\\
        \vspace{0.15cm}
        \textbf{Consecuencia:} el modelo usa pistas del futuro para «adivinar» el pasado.
      \end{block}
    \end{column}
    \begin{column}{0.48\textwidth}
      \begin{block}{Fuga por columna (incorrecto)}
        \small
        Predecir \texttt{consumo\_kwh} incluyendo como entrada \texttt{costo\_factura} (que se calculó multiplicando el consumo por la tarifa).\\
        \vspace{0.15cm}
        \textbf{Consecuencia:} error casi cero en papel, modelo inútil en la realidad.
      \end{block}
    \end{column}
  \end{columns}
  \vfill
  \centering
  \textbf{Pregunta de control:} ¿todas mis variables de entrada $X$ existen \emph{antes} del momento de predecir?
\end{frame}

\begin{frame}{Etapa 5: ningún modelo vale sin una referencia (\emph{Baseline})}
  \begin{columns}[T]
    \begin{column}{0.62\textwidth}
      \centering
      \begin{tikzpicture}
        \begin{axis}[
          ybar,
          bar width=28pt,
          width=\linewidth,
          height=5.2cm,
          ylabel={Error MAE en prueba (kWh)},
          symbolic x coords={Referencia (Media),Modelo Lineal},
          xtick=data,
          ymin=0,
          ymax=27,
          ymajorgrids=true,
          grid style={gris},
          nodes near coords,
          nodes near coords style={font=\small\bfseries},
          every node near coord/.append style={/pgf/number format/fixed, /pgf/number format/precision=2},
          tick label style={font=\scriptsize},
          label style={font=\scriptsize},
        ]
          \addplot[fill=naranja!75,draw=naranja] coordinates {
            (Referencia (Media),22.42)
            (Modelo Lineal,4.83)
          };
        \end{axis}
      \end{tikzpicture}
    \end{column}
    \begin{column}{0.35\textwidth}
      \small
      \textbf{Por qué exigir un baseline}
      \begin{itemize}
        \item Decir «mi modelo tiene MAE de 4.83» no significa nada por sí solo.
        \item Compararlo contra predecir la media del entrenamiento (MAE 22.42) demuestra una reducción del \textbf{78.5\%} del error en datos no vistos.
      \end{itemize}
    \end{column}
  \end{columns}
\end{frame}

\begin{frame}{Etapa 6: autopsia de resultados en datos de prueba}
  \begin{columns}[T]
    \begin{column}{0.58\textwidth}
      \centering
      \begin{tikzpicture}
        \begin{axis}[
          width=\linewidth,
          height=5.1cm,
          xlabel={Consumo real en prueba (kWh)},
          ylabel={Consumo predicho (kWh)},
          grid=major,
          grid style={gris},
          tick label style={font=\scriptsize},
          label style={font=\scriptsize},
          xmin=60, xmax=180,
          ymin=60, ymax=180,
        ]
          \addplot[domain=60:180, samples=2, dashed, color=verde, thick] {x};
          \addplot[
            only marks,
            mark=*,
            mark size=1.6pt,
            color=azul,
            opacity=0.8
          ] table[
            col sep=comma,
            x=consumo_kwh,
            y=prediccion,
          ] {predicciones_prueba.csv};
        \end{axis}
      \end{tikzpicture}
    \end{column}
    \begin{column}{0.38\textwidth}
      \small
      \textbf{Qué buscamos en esta gráfica}
      \begin{itemize}
        \item La línea punteada es la predicción perfecta ($\hat{y} = y$).
        \item Inspeccionamos los 3 puntos más alejados de la diagonal: ¿qué ocurrió en esas fechas?
        \item En la calificación final vale más explicar un fallo con rigor que esconderlo.
      \end{itemize}
    \end{column}
  \end{columns}
\end{frame}

\begin{frame}{Cronograma de 2 meses: del 10 de noviembre al 15 de enero}
  \begin{center}
    \scriptsize
    \begin{tabular}{llp{7.6cm}}
      \toprule
      \textbf{Hito} & \textbf{Plazo} & \textbf{Entregable verificable} \\
      \midrule
      \textbf{1. Banderazo y Datos} & Nov (Sem. 1--2) & Acta de proyecto (\texttt{plantilla\_propuesta.md}) + carga y auditoría del dataset crudo (Etapas 1 y 2). \\
      \addlinespace[3pt]
      \textbf{2. Limpieza y Referencia} & Fines Nov / Dic & Reglas de limpieza, análisis de atípicos, partición sin fuga y métrica del \emph{Baseline} (Etapas 3 y 4). \\
      \addlinespace[3pt]
      \textbf{3. Modelo e Iteración} & Diciembre & Entrenamiento del modelo principal, ajuste de hiperparámetros y comparación contra referencia (Etapa 5). \\
      \addlinespace[3pt]
      \textbf{4. Autopsia y Entrega} & \textbf{15 de enero} & Análisis de los peores errores en prueba, gráficas finales, límites del estudio y defensa (Etapa 6). \\
      \bottomrule
    \end{tabular}
  \end{center}
  \vspace{0.15cm}
  \begin{block}{Consejo práctico}
    \small Un pipeline sencillo que corre completo de la Etapa 1 a la 6 en noviembre es mucho mejor que un modelo gigante que no llega a evaluarse en enero.
  \end{block}
\end{frame}

\begin{frame}{¿Cómo se evaluará el proyecto el 15 de enero?}
  \begin{itemize}
    \item \textbf{Claridad de la pregunta y calidad de los datos (25\%):} auditoría explícita de faltantes, tipos, rangos y justificación de qué se conserva o descarta.
    \item \textbf{Diseño sin fuga y referencia justa (25\%):} separación correcta de entrenamiento/prueba y un \emph{baseline} bien calculado.
    \item \textbf{Implementación y comprensión del modelo (25\%):} código limpio, reproducible y capaz de explicarse línea por línea sin cajas negras mágicas.
    \item \textbf{Autopsia de errores y honestidad técnica (25\%):} inspección de fallos concretos y distinción clara entre capacidad predictiva y causalidad.
  \end{itemize}
\end{frame}

\begin{frame}{Taller de hoy: arrancar el motor}
  \begin{enumerate}
    \item Ejecutar el generador y el pipeline base de referencia:
      \begin{center}
        \texttt{python3 generar\_datos\_ejemplo.py} \quad y \quad \texttt{python3 pipeline\_base.py}
      \end{center}
    \item Inspeccionar las 6 etapas en consola y el archivo generado \texttt{predicciones\_prueba.csv}.
    \item Hacer un experimento rápido en \texttt{pipeline\_base.py}: quitar \texttt{equipos\_activos} de \texttt{ENTRADAS} y medir cuánto sube el MAE en prueba.
    \item Abrir \texttt{plantilla\_propuesta.md} y redactar el \textbf{Acta de Inicio} de su propio proyecto (Ruta A, B o C).
  \end{enumerate}
\end{frame}

\begin{frame}{Cierre: el banderazo inicial}
  \begin{center}
    \Large Pregunta clara $\rightarrow$ Datos auditados $\rightarrow$ Referencia $\rightarrow$ Modelo $\rightarrow$ Autopsia
  \end{center}
  \vfill
  \centering
  Hoy empieza la construcción de su proyecto final.\\
  Meta: llegar al \textbf{15 de enero} con un pipeline completo, reproducible y defendible.
\end{frame}

\end{document}
